\documentclass[11pt,a4paper,oneside]{scrreprt}

% Encoding, language, and typography
\usepackage{fontspec}
\usepackage[english]{babel}
\usepackage{microtype}
\usepackage{setspace}
\usepackage{parskip}
\setmainfont{TeX Gyre Pagella}
\setsansfont{TeX Gyre Heros}
\setmonofont{TeX Gyre Cursor}

% Page layout
\usepackage[a4paper,margin=25mm,headheight=15pt]{geometry}
\usepackage{fancyhdr}
\pagestyle{fancy}
\fancyhf{}
\lhead{\nouppercase{\leftmark}}
\rhead{Version 1.0}
\cfoot{\thepage}
\renewcommand{\headrulewidth}{0.4pt}
\renewcommand{\footrulewidth}{0pt}

% Tables, figures, and code
\usepackage{graphicx}
\usepackage{booktabs}
\usepackage{tabularx}
\usepackage{longtable}
\usepackage{array}
\usepackage{caption}
\usepackage{subcaption}
\usepackage{xcolor}
\usepackage{listings}
\lstset{
  basicstyle=\ttfamily\small,
  breaklines=true,
  frame=single,
  columns=fullflexible,
  showstringspaces=false
}

% Bibliography and references
\usepackage[
  backend=biber,
  style=authoryear,
  sorting=nyt
]{biblatex}
\addbibresource{references.bib}

\usepackage{hyperref}
\hypersetup{
  pdftitle={Professional Report Title},
  pdfauthor={Author or Organization},
  colorlinks=true,
  linkcolor=blue!45!black,
  citecolor=blue!45!black,
  urlcolor=blue!45!black
}
\usepackage[nameinlink,noabbrev]{cleveref}

% Index, glossary, and acronyms
\usepackage{imakeidx}
\makeindex
\usepackage[acronym,toc]{glossaries}
\makeglossaries
\newacronym{api}{API}{Application Programming Interface}

% Document metadata
\newcommand{\DocumentTitle}{Professional Report Title}
\newcommand{\DocumentSubtitle}{Subtitle or Project Name}
\newcommand{\DocumentAuthor}{Author or Organization}
\newcommand{\DocumentVersion}{1.0}
\newcommand{\DocumentStatus}{Draft}
\newcommand{\DocumentDate}{\today}
\newcommand{\DocumentClassification}{Internal}

\begin{document}

\begin{titlepage}
  \centering
  \vspace*{18mm}
  {\sffamily\bfseries\LARGE \DocumentTitle\par}
  \vspace{4mm}
  {\sffamily\Large \DocumentSubtitle\par}
  \vfill
  \begin{tabular}{ll}
    Author & \DocumentAuthor \\
    Version & \DocumentVersion \\
    Status & \DocumentStatus \\
    Classification & \DocumentClassification \\
    Date & \DocumentDate \\
  \end{tabular}
  \vfill
  {\small Prepared with LaTeX}
\end{titlepage}

\pagenumbering{roman}

\chapter*{Executive Summary}
\addcontentsline{toc}{chapter}{Executive Summary}
Summarize the purpose, scope, key findings, and recommendations in one page or less.

\tableofcontents
\listoffigures
\listoftables
\printglossary[type=\acronymtype,title={Acronyms}]
\printglossary

\clearpage
\pagenumbering{arabic}

\chapter{Introduction}
\label{chap:introduction}

Introduce the document purpose, audience, background, and scope. Define important terms such as \gls{api} and index important concepts\index{professional documentation}.

\section{Purpose}
\label{sec:purpose}

State what the document helps readers decide, understand, approve, or implement.

\section{Scope}
\label{sec:scope}

Describe what is included and excluded.

\chapter{Methodology}
\label{chap:methodology}

Explain the source material, analysis method, review approach, and assumptions.

\chapter{Findings}
\label{chap:findings}

Use structured headings, tables, figures, and references. See \cref{tab:summary} for a sample table and \cref{fig:example} for a sample figure placeholder.

\begin{table}[htbp]
  \centering
  \caption{Summary of Findings}
  \label{tab:summary}
  \begin{tabularx}{\textwidth}{l l X}
    \toprule
    ID & Priority & Finding \\
    \midrule
    F-001 & High & Replace this row with a concise finding and supporting evidence. \\
    F-002 & Medium & Add enough detail for readers to understand impact and next action. \\
    \bottomrule
  \end{tabularx}
\end{table}

\begin{figure}[htbp]
  \centering
  \fbox{\rule{0pt}{45mm}\rule{0.85\textwidth}{0pt}}
  \caption{Example figure placeholder}
  \label{fig:example}
\end{figure}

\chapter{Recommendations}
\label{chap:recommendations}

Prioritize recommendations by impact, effort, risk, and owner. Cite sources when factual claims depend on external evidence, for example \parencite{knuth1984texbook}.

\chapter{Conclusion}
\label{chap:conclusion}

Close with the most important decision, next step, or takeaway.

\printbibliography[heading=bibintoc]

\appendix

\chapter{Document Control}
\label{app:document-control}

\begin{longtable}{p{0.18\textwidth} p{0.18\textwidth} p{0.22\textwidth} p{0.32\textwidth}}
  \caption{Revision History}
  \label{tab:revision-history} \\
  \toprule
  Version & Date & Author & Notes \\
  \midrule
  \endfirsthead
  \toprule
  Version & Date & Author & Notes \\
  \midrule
  \endhead
  1.0 & \DocumentDate & \DocumentAuthor & Initial draft. \\
  \bottomrule
\end{longtable}

\printindex

\end{document}

